\begin{theorem}[canonical renormalization; Kahn I, \S7.4.3, p.~45]\label{thm:kahn-canonical-renormalization}
Let \((i,f)\colon(\BU^1,K')\to(\BU^0,K)\) be a \(\psi\)-quadratic-like
map admitting a quadratic-like restriction that is primitively
renormalizable of period \(p\); let \(K_p\) be the central small
filled Julia set, \(\KK_p\) the union of the small filled Julia sets,
and \(K_p'=i^{-1}(K_p)\) (iterate notation).  Then \(A=K_p\) and
\(B=\KK_p\) satisfy hypotheses (1)--(6) of the dynamical-localization
proposition for the iterate \((i^p,f^p)\colon\BU^p\to\BU^0\) of the
pullback tower, and the resulting dynamical localization
\((\widehat i,\widehat f)\colon(\BV',\widehat K_p')\to(\BV,\widehat K_p)\)
--- the \emph{canonical renormalization of \((i,f)\) of period \(p\)}
--- is a \(\psi\)-quadratic-like map.
\end{theorem}

\begin{theorem}[crucial length properties; Kahn I, \S7.4.3, p.~46]\label{thm:kahn-iterated-lengths}
Let \((i,f)\colon\BU'\to\BU\) be a \(\psi\)-quadratic-like map that is
primitively renormalizable of periods \(p\) and \(pq\) (\(p,q>1\)),
and let \((j,g)\colon\BV'\to\BV\) be its canonical renormalization of
period \(p\).  Then
\[
 |\gamma(g)|=|\gamma_p(f)|,
 \qquad
 |\gamma_q(g)|\ \ge\ |\gamma_{pq}(f)| .
\]
\end{theorem}

\begin{proof}[Proof (faithful paraphrase of Kahn I, \S7.4.3, p.~45 and 578--591)]
\emph{Canonical renormalization.}  Form the iterate
\((i^p,f^p)\colon\BU^p\to\BU^0\) via the pullback tower.  The data
\(A=K_p\), \(B=\KK_p\) satisfy the hypotheses of dynamical
localization for this iterate: \(K_p\) is an FJ-set and a component of
the compact set \(\KK_p\) (by primitivity the small filled Julia sets
are pairwise disjoint); \(\BU^0\sm\KK_p\) is connected; \(K_p'\) is a
component of \((f^p)^{-1}(K_p)\); and the branch values of \(f^p\) lie
in \(\KK_p\).  The dynamical localization of \(f^p\) at \(K_p\) with
removal of \(\KK_p\) is then \(\psi\)-polynomial-like of degree equal
to the local degree of \(f^p\) at \(K_p'\), namely \(2\); this is the
canonical renormalization (\S7.4.3, pp.~45--46).

\emph{Length properties.}  The equality \(|\gamma(g)|=|\gamma_p(f)|\)
is exactly conclusion (c) of the dynamical-localization proposition
applied to this localization (\S7.4.3, p.~46).  For the inequality, write
\(K(g),\KK_q(g)\) for the objects of \(g\) and
\(K,\KK_p,\KK_{pq}\) for those of \(f\), and let
\[
 \pi\colon(\BV,K(g))\longrightarrow\bigl((\BU\sm\KK_p)\cup K_p,\ K_p\bigr)
\]
be the local cover constructed in the localization (\S7.4.3, p.~46).
Then \(\pi(\BV\sm\KK_q(g))\subseteq\BU\sm\KK_{pq}\): indeed
\(\pi(\BV\sm K(g))\subseteq\BU\sm\KK_p\subseteq\BU\sm\KK_{pq}\),
while \(\pi(K(g)\sm\KK_q(g))=K_p\sm\KK_{pq}\) (\S7.4.3, p.~46).  Thus
\(\pi\) maps the hyperbolic surface \(\BV\sm\KK_q(g)\)
holomorphically into \(\BU\sm\KK_{pq}\), carrying \(\gamma_q(g)\) into
the free homotopy class of \(\gamma_{pq}(f)\); the Schwarz--Pick lemma
(Poincar\'e-metric contraction under holomorphic maps) gives
\(|\gamma_q(g)|\ge|\gamma_{pq}(f)|\) (\S7.4.3, p.~46).
\end{proof}